\section{Conclusion}
\label{sec:conclude}

This paper shows that production coding agents can be understood as answers to a recurring set of design questions: where reasoning sits relative to the harness, how execution, safety, extensibility, context, delegation, and persistence are organized, and which trade-offs those choices encode. Claude Code occupies a clear design point within that space. It gives the model broad local autonomy while surrounding it with a dense deterministic harness for permissioning, tool routing, context compaction, extensibility, and session recovery. Read through the five values and thirteen design principles identified in \Cref{sec:values}, these choices are coherent rather than ad hoc: the system consistently prioritizes human decision authority, safety, reliable execution, capability amplification, and contextual adaptability.

The OpenClaw and Hermes Agent comparisons show that the same design questions recur in different agent systems but produce different answers. Where Claude Code invests in per-action safety classification and graduated context compression within a CLI harness, OpenClaw invests in perimeter-level access control and structured long-term memory within a multi-channel gateway, and Hermes Agent renders per-action approvals across many surfaces, from one process, with pluggable memory and model backends. The three systems compose through ACP at multiple positions: OpenClaw can host Claude Code as an external harness via ACP, and Hermes Agent sits on both sides of the host/guest split. For agent builders, the most consequential open question is therefore not how to add more autonomy, but how to preserve human capability over time. Taken together, the long-term capability question introduced in \Cref{sec:values:lens}, the analysis in \Cref{sec:discuss}, and the open questions surveyed in \Cref{sec:future} show that the architecture provides limited mechanisms that explicitly preserve long-term human understanding, codebase coherence, or the developer pipeline. Future systems could treat that sustainability gap as a first-class design problem, not a downstream evaluation metric.
